\section*{Hoofdstuk \ref{ch:b2:euler-bernoulli} --- Volmaakte fluïda: Euler
en Bernoulli}

\begin{solution}{exo:b2:euler-bernoulli:1}
$v_2 = 1.5 \times 4 = \qty{6.0}{m/s}$,
$P_2 = 3.0 \times 10^5 - 500(36 - 2.25) =
\qty{2.83}{bar}$. Bij \qty{2}{cm}:
$v = \qty{24}{m/s}$,
$P = 3.0 \times 10^5 - 500
(576 - 2.25) = \qty{0.13}{bar}$ --- dicht bij de
dampdruk. Een negatieve uitkomst betekent dat het veronderstelde debiet
onmogelijk is: het water caviteert en de stroming wordt afgeknepen.
\end{solution}

\begin{solution}{exo:b2:euler-bernoulli:2}
$v = \sqrt{2 \times 15000/0.74} = \qty{201}{m/s}$; aangewezen
$\sqrt{2 \times 15000/
1.22} = \qty{157}{m/s}$. De aangewezen luchtsnelheid
meet de \omterm{thm:b2:euler-bernoulli:bernoulli}{dynamische druk}, en dat is wat de vleugel voelt: de overtreksnelheid
is op elke hoogte dezelfde \emph{aangewezen} snelheid.
\end{solution}

\begin{solution}{exo:b2:euler-bernoulli:3}
Diepte \qty{2.0}{m}: $v = \sqrt{2g \times 2} = \qty{6.3}{m/s}$;
$D_V = 0.6 \times 10^{-4}
\times 6.3 = \qty{0.38}{L/s}$; valtijd
$\sqrt{2/g} = \qty{0.45}{s}$, worp \qty{2.8}{m}. De worp
$= 2\sqrt{h(H - h)}$ is symmetrisch in $h \leftrightarrow H -
h$: een gat op
diepte \qty{1.0}{m} (\qty{2}{m} boven de vloer) treft hetzelfde punt.
\end{solution}

\begin{solution}{exo:b2:euler-bernoulli:4}
$\Delta P = \tfrac12 \times 1.2 \times 900 = \qty{540}{Pa}$, met de laagste
druk buiten: netto kracht \qty{54}{kN} \emph{omhoog}; het dak weegt
\qty{49}{kN}: het licht op. (Echte daken bezwijken eerst aan hun randen,
waar de stroming loslaat.)
\end{solution}

\begin{solution}{exo:b2:euler-bernoulli:5}
(a) $\tan\theta = 3/9.81$, $\theta = \qty{17}{\degree}$;
$2 \times 0.306 = \qty{0.61}{m}$ (achteraan hoger). (b)
$\Delta P = \rho aL = \qty{6.0}{kPa}$. (c) $\tan\theta =
0.61$, verschil
\qty{1.22}{m}: de voorkant stijgt \qty{0.61}{m} boven het gemiddelde van
\qty{0.40}{m}: \qty{1.01}{m} $>$ \qty{0.8}{m}, het water raakt het deksel.
(d) $a = v^2/R =
\qty{5.6}{m/s^2}$: $\theta = \qty{30}{\degree}$, de
buitenbochtzijde hoger.
\end{solution}

\begin{solution}{exo:b2:euler-bernoulli:6}
(a) Paraboloïde $z = z_0 + \Omega^2r^2/2g$,
$\Omega = 4\pi = \qty{12.6}{rad/s}$. (b)
$\Omega^2R^2/2g = 158 \times 0.0225/19.6 = \qty{0.18}{m}$. (c) De gemiddelde
hoogte van een paraboloïde over de schijf ligt halverwege: middelpunt
$0.20 - 0.09 = \qty{0.11}{m}$, rand \qty{0.29}{m}. (d) Het midden valt droog
wanneer $\Omega^2R^2/4g = h_0$:
$\Omega = \sqrt{4gh_0}/R =
\qty{18.7}{rad/s} = 3.0$ omwentelingen per
seconde. (e) Een paraboloïde bundelt evenwijdige stralen in één punt, op
$f = g/2\Omega^2$; vloeibarespiegeltelescopen laten kwik draaien.
\end{solution}

\begin{solution}{exo:b2:euler-bernoulli:7}
(a) $v = \sqrt{2g \times 3} = \qty{7.7}{m/s}$;
$D_V = \pi \times 10^{-4} \times 7.7 =
\qty{2.4}{L/s}$. (b)
$P = P_0 - \rho g(2) - \tfrac12\rho v^2 = 100 - 19.6 - 29.4 =
\qty{51}{kPa}$.
(c) $P_0 - \rho gh_c - \tfrac12\rho v^2 \ge \qty{2.3}{kPa}$:
$h_c \le
\qty{7.0}{m}$. (d) Lucht is samendrukbaar en licht: er is geen
doorlopende vloeistofkolom om de druk over te brengen, en het water aan
beide zijden blijft gewoon staan.
\end{solution}

\begin{solution}{exo:b2:euler-bernoulli:8}
(a) $\Delta P = mg/S = \qty{654}{Pa}$; $\tfrac12\rho v^2 = \qty{2.2}{kPa}$,
$C_L = 0.30$. (b) $v_{\text{boven}}^2 = 58^2 + 2 \times 654/1.22$:
\qty{66.6}{m/s}. (c) $F_L' =
\qty{981}{N/m} = \rho v\Gamma$:
$\Gamma = \qty{13.4}{m^2/s}$. (d)
$v = \sqrt{2mg/
\rho SC_L} = \sqrt{19620/25.6} = \qty{28}{m/s}$.
\end{solution}

\begin{solution}{exo:b2:euler-bernoulli:9}
(a) $v_2 = \qty{2.5}{m/s}$;
$\Delta P = \tfrac12 \times 1060 \times (6.25 - 0.25) =
\qty{3.2}{kPa}$,
$P_2 = \qty{9.8}{kPa}$. (b) Nog altijd boven \qty{1}{kPa}: nee. (c) Bij
\qty{0.05}{cm^2} is $v = \qty{5}{m/s}$ en $\Delta P = \qty{13}{kPa}$:
$P_2 \approx 0$, onder de weefseldruk --- de wand klapt dicht, de stroming
stopt, de druk herstelt zich, de wand gaat weer open: een gefladder (het
geruis dat een arts hoort).
\end{solution}

\begin{solution}{exo:b2:euler-bernoulli:10}
(a) $c_PT_0 = c_PT + v^2/2$ met $c_P = \gamma R/(\gamma - 1)M_{\text{mol}}$
en $c^2 = \gamma RT/M_{\text{mol}}$: $v^2/2c_PT = \tfrac12(\gamma - 1)M^2$.
(b) $M =
0.85$: $T_0 = 220 \times 1.145 = \qty{252}{K}$; $M = 2$:
$220 \times 1.8 = \qty{396}{K}$; $M = 25$: $220 \times 126 = \qty{28000}{K}$
--- zinledig: de lucht dissocieert en ioniseert, $c_P$ is niet constant
(echte stuwtemperaturen bedragen enkele duizenden kelvin). (c)
$P_0/P = (1 + \tfrac{\gamma - 1}2M^2)^{\gamma/(\gamma
- 1)} \approx 1 + \tfrac{\gamma}2M^2 + \tfrac{\gamma}8M^4$,
en $\gamma PM^2 = \rho v^2$: $P_0 - P = \tfrac12\rho v^2(1 + M^2/4)$. Leest
men $v$ af uit $\sqrt{2(P_0 - P)/\rho}$, dan overschat men haar met
$\sqrt{1 + M^2/4} - 1 \approx 9\%$ bij $M = 0.85$.
\end{solution}

\begin{solution}{exo:b2:euler-bernoulli:11}
(a) Van het oppervlak ($v \approx 0$, $P_0$, hoogte $h$) tot de uitmonding
($v$, $P_0$, hoogte $0$), met $\int\partial_tv\,\dd s = \ell\,\dd v/\dd t$.
(b) $\dd v/(V^2 -
v^2) = \dd t/2\ell$, $V^2 = 2gh$:
$\operatorname{artanh}(v/V) = Vt/2\ell$. (c) $V =
\qty{9.9}{m/s}$;
$\tanh x = 0.9$ bij $x = 1.47$: $t = 2\ell x/V = \qty{15}{s}$. (d)
$\dd v/\dd t \approx \qty{-20}{m/s^2}$:
$\Delta P = \rho\ell \times 20 = \qty{9.9}{bar}$. Water is licht
samendrukbaar: een snelle sluiting stuurt een drukgolf op pad,
$\Delta P = \rho c\Delta v \approx \qty{140}{bar}$.
\end{solution}

\begin{solution}{exo:b2:euler-bernoulli:12}
(a) Stationair, volmaakt, onsamendrukbaar en rotatievrij: één constante voor
de hele stroming. (b) Op het oppervlak is $P = P_0$:
$\tfrac12\rho v^2 + \rho gz =
0$, $z = -\Gamma^2/8\pi^2gr^2$. (c) Radiale
Euler: $\partial_rP = \rho\Omega^2r$, verticaal hydrostatisch: het oppervlak
voldoet aan $z(r) - z(a) = \Omega^2(r^2 -
a^2)/2g$. (d) Bad:
$\Omega = \Gamma/2\pi a^2 = \qty{191}{rad/s}$;
$z(a) = -\Omega^2a^2/
2g = \qty{-4.7}{cm}$, en de kern voegt nog eens
\qty{4.7}{cm} toe: \qty{9.3}{cm} diep. Tornado:
$\Delta P = \tfrac12\rho v_{\max}^2$ buiten plus evenveel binnen:
$\rho v_{\max}^2 = 1.2 \times 6400 = \qty{7.7}{kPa}$ onder de omgevingsdruk
in het middelpunt.
\end{solution}

\begin{solution}{pb:b2:euler-bernoulli:1}
\textbf{1.} $\rho g \times 150 = \qty{14.7}{bar}$ boven de atmosferische
druk; $\tfrac12\rho gH_d^2
\times \qty{1}{m} = \qty{1.1e8}{N}$.

\textbf{2.} $\rho gH = \qty{19.6}{bar}$ overdruk, \qty{20.6}{bar} absoluut.

\textbf{3.} $v = 60/7.07 = \qty{8.5}{m/s}$;
$\tfrac12\rho v^2 = \qty{0.36}{bar}$;
$P =
P_0 + \rho gH - \tfrac12\rho v^2 = \qty{20.3}{bar}$ absoluut.

\textbf{4.} $v_j = \sqrt{2gH} = \qty{62.6}{m/s}$;
$s = 60/62.6 = \qty{0.96}{m^2}$, $d =
\qty{1.1}{m}$.

\textbf{5.}
$\tfrac12\rho D_Vv_j^2 = \tfrac12\rho D_V(2gH) = \rho gD_VH = \qty{118}{MW}$.

\textbf{6.} $m = \rho S\ell = \qty{2.8e6}{kg}$,
$E_k = \tfrac12mv^2 = \qty{1.0e8}{J}$; doorlooptijd $\ell/v = \qty{47}{s}$.

\textbf{7.} $\rho gH = (P - P_0) + \tfrac12\rho v^2$: \qty{19.3}{bar} aan
druk en \qty{0.36}{bar} aan kinetische energie; in het mondstuk daalt de
druk tot $P_0$ terwijl $\tfrac12\rho v^2$ tot $\rho gH$ stijgt: vlak erna is
$P = P_0$.

\textbf{8.}
$\Delta h = \lambda(\ell/D)v^2/2g = 0.015 \times 133 \times 3.67 = \qty{7.3}{m}$,
$3.7\%$ van het vermogen.

\textbf{9.} $v = \qty{19.1}{m/s}$, $v^2/2g = \qty{18.6}{m}$,
$\Delta h = 0.015 \times 200 \times
18.6 = \qty{56}{m}$: $28\%$ verloren ---
wat aan staal wordt bespaard, wordt elke seconde aan energie betaald.

\textbf{10.} $v_t = 60/\pi = \qty{19.1}{m/s}$;
$\Delta P = 500(19.1^2 - 8.49^2) =
\qty{146}{kPa}$.

\textbf{11.} $\Delta P = (\rho_{\text{Hg}} - \rho)g\,\Delta h$:
$\Delta h = 146000/(12600
\times 9.81) = \qty{1.18}{m}$.

\textbf{12.} $\Delta P \propto v^2 \propto D_V^2$:
$D_V \propto \sqrt{\Delta h}$; het halve debiet geeft een kwart,
\qty{0.30}{m}.

\textbf{13.} Een uitstekende aftakking staat in de stroming of verstoort
haar en leest een deel van de \omterm{thm:b2:euler-bernoulli:bernoulli}{dynamische druk} af (een \omterm{prop:b2:euler-bernoulli:pitot}{pitotbuis}), niet de
statische druk.

\textbf{14.}
$P = P_0 - \rho g \times 6 - \tfrac12\rho v^2 = 100 - 59 - 36 = \qty{5}{kPa}$:
vlak bij de dampdruk, opgeloste lucht komt vrij en het water kan gaan koken
(cavitatie), waardoor de kolom breekt. Verlaag de top, of plaats een
ontluchter met een vacuümpomp, of verminder het debiet.

\textbf{15.} $D_V/A = 60/2 \times 10^6 = \qty{3e-5}{m/s} = \qty{11}{cm/h}$.

\textbf{16.} $A\,\dd h/\dd t = -s\sqrt{2gh}$:
$\sqrt h = \sqrt{h_0} - (s/A)\sqrt{g/2}\,t$.

\textbf{17.}
$\Delta t = (A/s)\sqrt{2/g}(\sqrt{200} - \sqrt{190}) = 2.09 \times 10^6 \times
0.45 \times 0.36 = \qty{3.4e5}{s} = \qty{3.9}{dagen}$;
bij constant ontwerpdebiet $10A/D_V = \qty{3.3e5}{s}$: vrijwel hetzelfde,
omdat het mondstuk juist op de stroming van Torricelli bij $h \approx H$ is
gedimensioneerd.

\textbf{18.} $\ell\,\dd v/\dd t = gH - v^2/2$: $v = V\tanh(Vt/2\ell)$,
$V = \qty{62.6}{m/s}$, tijdconstante $2\ell/V = \qty{13}{s}$.

\textbf{19.} Bij $v = V/2$:
$\dd v/\dd t = (gH - V^2/8)/\ell = \qty{3.7}{m/s^2}$;
$P = P_0 + \rho gH - \tfrac12\rho v^2 - \rho\ell\,\dd v/\dd t = 1.0 + 19.6 - 4.9 - 14.7
= \qty{1.0}{bar}$:
het grootste deel van de vervalhoogte is bezig de kolom te versnellen, zodat
de druk onderaan ver onder haar stationaire waarde ligt.

\textbf{20.} $|\dd v/\dd t| = 8.49/10 = \qty{0.85}{m/s^2}$:
$\Delta P = \rho\ell \times 0.85 =
\qty{3.4}{bar}$.

\textbf{21.} $\Delta P = 1000 \times 1400 \times 8.49 = \qty{119}{bar}$, zes
keer de statische druk. De heen-en-terugreis van de golf duurt
$2\ell/c = \qty{0.57}{s}$: een sluiting die sneller is dan dat, is
``plotseling'' en krijgt de volle overdruk van Joukowsky; tragere sluitingen
krijgen minder.

\textbf{22.} $Av = A_s\,\dd z/\dd t$.

\textbf{23.} Euler langs de tunnel, van het stuwmeer (druk
$P_0 + \rho g\times$diepte) tot de voet van de schacht (druk
$P_0 + \rho g
(\text{diepte} + z)$): $\rho L\,\dd v/\dd t = -\rho gz$; met
$v = (A_s/A)\dot z$: $\ddot z + (gA/LA_s)z = 0$;
$T = 2\pi\sqrt{LA_s/gA} = 2\pi\sqrt{2000 \times 100/
68.7} = \qty{340}{s} \approx \qty{5.7}{min}$.

\textbf{24.} $\tfrac12\rho ALv^2 = \tfrac12\rho gA_sz_{\max}^2$ met
$v = 60/7 =
\qty{8.6}{m/s}$: $z_{\max} = v\sqrt{AL/gA_s} = \qty{32}{m}$.

\textbf{25.} Dam: hydrostatica, $\rho gh$. Turbine met gesloten afsluiters:
hetzelfde, $\rho gH$. Drukleiding in bedrijf: Bernoulli, $\tfrac12\rho v^2$
genomen uit de statische vervalhoogte. Waterslag: samendrukbaarheid,
$\rho c\Delta v$. Waterslot: niet-stationaire Euler langs de tunnel, een
oscillatie met periode $2\pi\sqrt{LA_s/gA}$.
\end{solution}
